% Options for packages loaded elsewhere
\PassOptionsToPackage{unicode}{hyperref}
\PassOptionsToPackage{hyphens}{url}
\PassOptionsToPackage{dvipsnames,svgnames,x11names}{xcolor}
%
\documentclass[
  10pt,
]{article}
\usepackage{amsmath,amssymb}
\usepackage{iftex}
\ifPDFTeX
  \usepackage[T1]{fontenc}
  \usepackage[utf8]{inputenc}
  \usepackage{textcomp} % provide euro and other symbols
\else % if luatex or xetex
  \usepackage{unicode-math} % this also loads fontspec
  \defaultfontfeatures{Scale=MatchLowercase}
  \defaultfontfeatures[\rmfamily]{Ligatures=TeX,Scale=1}
\fi
\usepackage{lmodern}
\ifPDFTeX\else
  % xetex/luatex font selection
  \setmainfont[]{DejaVu Serif}
  \setsansfont[]{DejaVu Sans}
  \setmonofont[]{DejaVu Sans Mono}
  \setmathfont[]{Latin Modern Math}
\fi
% Use upquote if available, for straight quotes in verbatim environments
\IfFileExists{upquote.sty}{\usepackage{upquote}}{}
\IfFileExists{microtype.sty}{% use microtype if available
  \usepackage[]{microtype}
  \UseMicrotypeSet[protrusion]{basicmath} % disable protrusion for tt fonts
}{}
\makeatletter
\@ifundefined{KOMAClassName}{% if non-KOMA class
  \IfFileExists{parskip.sty}{%
    \usepackage{parskip}
  }{% else
    \setlength{\parindent}{0pt}
    \setlength{\parskip}{6pt plus 2pt minus 1pt}}
}{% if KOMA class
  \KOMAoptions{parskip=half}}
\makeatother
\usepackage{xcolor}
\usepackage[margin=2.0cm]{geometry}
\usepackage{color}
\usepackage{fancyvrb}
\newcommand{\VerbBar}{|}
\newcommand{\VERB}{\Verb[commandchars=\\\{\}]}
\DefineVerbatimEnvironment{Highlighting}{Verbatim}{commandchars=\\\{\}}
% Add ',fontsize=\small' for more characters per line
\newenvironment{Shaded}{}{}
\newcommand{\AlertTok}[1]{\textcolor[rgb]{1.00,0.00,0.00}{\textbf{#1}}}
\newcommand{\AnnotationTok}[1]{\textcolor[rgb]{0.38,0.63,0.69}{\textbf{\textit{#1}}}}
\newcommand{\AttributeTok}[1]{\textcolor[rgb]{0.49,0.56,0.16}{#1}}
\newcommand{\BaseNTok}[1]{\textcolor[rgb]{0.25,0.63,0.44}{#1}}
\newcommand{\BuiltInTok}[1]{\textcolor[rgb]{0.00,0.50,0.00}{#1}}
\newcommand{\CharTok}[1]{\textcolor[rgb]{0.25,0.44,0.63}{#1}}
\newcommand{\CommentTok}[1]{\textcolor[rgb]{0.38,0.63,0.69}{\textit{#1}}}
\newcommand{\CommentVarTok}[1]{\textcolor[rgb]{0.38,0.63,0.69}{\textbf{\textit{#1}}}}
\newcommand{\ConstantTok}[1]{\textcolor[rgb]{0.53,0.00,0.00}{#1}}
\newcommand{\ControlFlowTok}[1]{\textcolor[rgb]{0.00,0.44,0.13}{\textbf{#1}}}
\newcommand{\DataTypeTok}[1]{\textcolor[rgb]{0.56,0.13,0.00}{#1}}
\newcommand{\DecValTok}[1]{\textcolor[rgb]{0.25,0.63,0.44}{#1}}
\newcommand{\DocumentationTok}[1]{\textcolor[rgb]{0.73,0.13,0.13}{\textit{#1}}}
\newcommand{\ErrorTok}[1]{\textcolor[rgb]{1.00,0.00,0.00}{\textbf{#1}}}
\newcommand{\ExtensionTok}[1]{#1}
\newcommand{\FloatTok}[1]{\textcolor[rgb]{0.25,0.63,0.44}{#1}}
\newcommand{\FunctionTok}[1]{\textcolor[rgb]{0.02,0.16,0.49}{#1}}
\newcommand{\ImportTok}[1]{\textcolor[rgb]{0.00,0.50,0.00}{\textbf{#1}}}
\newcommand{\InformationTok}[1]{\textcolor[rgb]{0.38,0.63,0.69}{\textbf{\textit{#1}}}}
\newcommand{\KeywordTok}[1]{\textcolor[rgb]{0.00,0.44,0.13}{\textbf{#1}}}
\newcommand{\NormalTok}[1]{#1}
\newcommand{\OperatorTok}[1]{\textcolor[rgb]{0.40,0.40,0.40}{#1}}
\newcommand{\OtherTok}[1]{\textcolor[rgb]{0.00,0.44,0.13}{#1}}
\newcommand{\PreprocessorTok}[1]{\textcolor[rgb]{0.74,0.48,0.00}{#1}}
\newcommand{\RegionMarkerTok}[1]{#1}
\newcommand{\SpecialCharTok}[1]{\textcolor[rgb]{0.25,0.44,0.63}{#1}}
\newcommand{\SpecialStringTok}[1]{\textcolor[rgb]{0.73,0.40,0.53}{#1}}
\newcommand{\StringTok}[1]{\textcolor[rgb]{0.25,0.44,0.63}{#1}}
\newcommand{\VariableTok}[1]{\textcolor[rgb]{0.10,0.09,0.49}{#1}}
\newcommand{\VerbatimStringTok}[1]{\textcolor[rgb]{0.25,0.44,0.63}{#1}}
\newcommand{\WarningTok}[1]{\textcolor[rgb]{0.38,0.63,0.69}{\textbf{\textit{#1}}}}
\usepackage{longtable,booktabs,array}
\usepackage{calc} % for calculating minipage widths
% Correct order of tables after \paragraph or \subparagraph
\usepackage{etoolbox}
\makeatletter
\patchcmd\longtable{\par}{\if@noskipsec\mbox{}\fi\par}{}{}
\makeatother
% Allow footnotes in longtable head/foot
\IfFileExists{footnotehyper.sty}{\usepackage{footnotehyper}}{\usepackage{footnote}}
\makesavenoteenv{longtable}
\setlength{\emergencystretch}{3em} % prevent overfull lines
\providecommand{\tightlist}{%
  \setlength{\itemsep}{0pt}\setlength{\parskip}{0pt}}
\setcounter{secnumdepth}{-\maxdimen} % remove section numbering
% doc_header.tex -- LaTeX preamble additions used by build_docs.sh for every
% Markdown -> LaTeX -> PDF conversion (the documents of the folder Revision_SM_Higgs).
\usepackage{etoolbox}
% longtable may not be loaded by pandoc when a document has no table
\usepackage{longtable}
% bold vectors: \bm -> \symbf under unicode-math (XeLaTeX), bm package otherwise
\makeatletter
\@ifpackageloaded{unicode-math}{\providecommand{\bm}[1]{\symbf{#1}}}{\usepackage{bm}}
\makeatother
% smaller type in long tables so that wide audit tables fit the text width;
% inside a table a line break is allowed after an underscore of plain text
% (pandoc writes it as \_), so that long identifiers wrap inside their column
% rows are set a little apart, so that symbols with indices do not touch
\AtBeginEnvironment{longtable}{\footnotesize\renewcommand{\arraystretch}{1.45}%
  \let\sptunderscore\_\def\_{\sptunderscore\allowbreak}}
\setlength{\LTcapwidth}{\textwidth}
% A short table is kept on one page.  The build (code/scripts/keep_tables.py) ends the rows of such a table with \\*,
% which forbids a page break after a row, and sets \sptkeptpenalty to 10000 in front of the table.  This penalty is
% put here in the place of the breakpoint (penalty 0) that longtable leaves after every table, so that the paragraph
% that follows a kept table, its caption, begins on the same page; it is set back to 0 after the table.
\newcount\sptkeptpenalty
\sptkeptpenalty=0
\makeatletter
\patchcmd{\endlongtable}{\endgraf\penalty\z@\addvspace\LTpost}%
  {\endgraf\penalty\sptkeptpenalty\addvspace\LTpost\global\sptkeptpenalty=\z@}%
  {}{\PackageError{doc_header}{the end of longtable could not be patched}{The tables can still be typeset, but a short table may be separated from its caption.}}
\makeatother
% allow line breaks after slashes and hyphens inside table cells
\usepackage{xurl}
% code blocks: break over-long lines instead of running into the margin.
% Plain ``` blocks are typeset by pandoc as {verbatim}, highlighted ones
% (```bash) as {Highlighting}; both are re-declared as breakable fancyvrb
% environments.  The continuation line is marked with a small hooked arrow.
\usepackage{fvextra}
\RecustomVerbatimEnvironment{verbatim}{Verbatim}{breaklines,breakanywhere,fontsize=\small}
\makeatletter
\@ifundefined{Highlighting}{}{%
  \RecustomVerbatimEnvironment{Highlighting}{Verbatim}{breaklines,breakanywhere,fontsize=\small,commandchars=\\\{\}}}
\makeatother
% inline code (pandoc writes it as \texttt{...}): never hyphenate -- an inserted
% hyphen could be read as part of a file name -- but allow a line break after
% an underscore, slash, comma, colon or equals sign, so that long paths and
% column lists wrap inside the text width and inside table columns.
% Inside a formula (e.g. the name of a notebook symbol in an equation) the same
% command must first leave math mode, which \text of amsmath does.
\ExplSyntaxOn
\tl_new:N \l__spt_code_tl
\cs_new_protected:Npn \__spt_texttt:n #1
  {
    \tl_set:Nn \l__spt_code_tl {#1}
    \regex_replace_all:nnN { ([\/\,\:\=]) } { \1 \c{allowbreak} } \l__spt_code_tl
    \regex_replace_all:nnN { \c{_} } { \c{_} \c{allowbreak} } \l__spt_code_tl
    \group_begin: \ttfamily \hyphenchar \font = -1 \scan_stop: \l__spt_code_tl \group_end:
  }
\RenewDocumentCommand \texttt { m }
  {
    \mode_if_math:TF { \text { \__spt_texttt:n {#1} } } { \__spt_texttt:n {#1} }
  }
\ExplSyntaxOff
% A figure is written in the Markdown source as an image paragraph followed by
% a caption paragraph ("**Figure S1.** ..."), so that the caption is also
% visible on GitHub.  An image that starts a paragraph is centered.
\usepackage{graphicx}% pandoc loads it only for documents that contain an image
\NewCommandCopy\sptincludegraphics\includegraphics
\RenewDocumentCommand\includegraphics{O{} m}{%
  \ifvmode\begin{center}\sptincludegraphics[#1]{#2}\end{center}%
  \else\sptincludegraphics[#1]{#2}\fi}
% keep figures where they are referenced
\usepackage{float}
\floatplacement{figure}{H}
% running header with the document title
\usepackage{fancyhdr}
\pagestyle{fancy}
\fancyhf{}
\fancyhead[L]{\small PROJECT spacetimes, Revision\_SM\_Higgs}
\fancyhead[R]{\small\thepage}
\renewcommand{\headrulewidth}{0.4pt}
% 16x16 matrices: tighter columns so that a full matrix fits the text width
\setlength{\arraycolsep}{3.2pt}
% The documents of this project contain long tagged symbols, file paths and
% identifiers that cannot be hyphenated.  A paragraph is therefore allowed to
% have looser lines before TeX gives up and lets a line stick out of the margin.
\AtBeginDocument{\tolerance=2000 \emergencystretch=5em \hbadness=2000}
\ifLuaTeX
  \usepackage{selnolig}  % disable illegal ligatures
\fi
\IfFileExists{bookmark.sty}{\usepackage{bookmark}}{\usepackage{hyperref}}
\IfFileExists{xurl.sty}{\usepackage{xurl}}{} % add URL line breaks if available
\urlstyle{same}
\hypersetup{
  pdftitle={Revision\_SM\_Higgs: the superfluid of the spinor fields on the 3+1 manifold of the Lorentz super-loop, revised so that it has the quantum numbers of the Higgs doublet of the Standard Model},
  pdfauthor={Prepared for Patrick Lee Nash, Ph.D.},
  colorlinks=true,
  linkcolor={blue},
  filecolor={Maroon},
  citecolor={blue},
  urlcolor={blue},
  pdfcreator={LaTeX via pandoc}}

\title{Revision\_SM\_Higgs: the superfluid of the spinor fields on the
3+1 manifold of the Lorentz super-loop, revised so that it has the
quantum numbers of the Higgs doublet of the Standard Model}
\usepackage{etoolbox}
\makeatletter
\providecommand{\subtitle}[1]{% add subtitle to \maketitle
  \apptocmd{\@title}{\par {\large #1 \par}}{}{}
}
\makeatother
\subtitle{README of the folder Revision\_SM\_Higgs of PROJECT
spacetimes}
\author{Prepared for Patrick Lee Nash, Ph.D.}
\date{November 21, 2026}

\begin{document}
\maketitle

\hypertarget{the-revised-hypotheses-read-this-first}{%
\section{0. The revised hypotheses (read this
first)}\label{the-revised-hypotheses-read-this-first}}

This folder was made on 21 November 2026 as a copy of the folder
\texttt{Revision\_HS\_CJT\_and\_Kohn-Sham} at the commit
\texttt{19db425}, which stays unchanged as the record of the earlier
hypotheses; the commit also carries the git tag
\texttt{backup/pre-SM-Higgs-2026-11-21}. The author asked that the
hypotheses be modified so that the superfluid can be the Higgs field of
the Standard Model, with the same quantum numbers as the Higgs, and
chose the construction:

\begin{itemize}
\tightlist
\item
  the isodoublet \(\Psi_{\rm LL}\) is joined by a singlet field
  \(\chi_{\rm SH}\), both Dirac fields and colour singlets;
\item
  \(SU(2)_{\rm v}\) is \(SU(2)_L\), and the fermion number is the
  hypercharge, \(Y^{\rm SM}=F_{\rm LL}/4\), so both fields have
  \(Y^{\rm SM}=1/4\) and the electric charges are \(3/4\), \(-1/4\) (the
  doublet) and \(1/4\) (the singlet);
\item
  the superfluid is the condensate of the spin-zero pairs of a particle
  of \(\Psi_{\rm LL}\) and a particle of \(\chi_{\rm SH}\): its pair
  field \(\Phi_{\rm SH}\) is a Lorentz scalar, a colour singlet and an
  \(SU(2)_L\) doublet with \(Y^{\rm SM}=+1/2\), the quantum numbers of
  the Higgs doublet;
\item
  the gauge couplings are the measured ones, and the scale of the model
  is fixed by \(v_{\rm SH}=246.22\) GeV.
\end{itemize}

The hypotheses, their reasons and their immediate consequences are
stated in \texttt{provenance/H0\_hypotheses.md} (\texttt{.tex},
\texttt{.pdf}); the request, the choices and the plan are in the task
record \texttt{provenance/P0\_task\_record.md}.

\textbf{What is current and what is inherited.} The hypotheses document
and the task record are current. The other documents, the programs, the
data, the notebook and the article in \texttt{paper\_prd/} are inherited
unchanged from the earlier hypotheses and are rewritten in stages: (1)
the folder, the hypotheses and the task record; (2) the quantum numbers
and the pattern of the breaking, checked by programs in Rust, Python and
the Wolfram Language; (3) the mean-field superfluid with the methods HS,
KS and CJT; (4) the effective potential, the mass of the amplitude mode
against \(m_H\), the masses of \(W^\pm_\mu\) and \(Z_\mu\), the massless
photon and the parameter \(\rho\) (the former stages 4 and 5, merged);
(6) the dynamics in time, the notebook and the article. Every inherited
document begins with a note that says so. Stages 1 to 4 and stage 6 (the
dynamics in time, the notebook and the article) are complete for the
second version of the hypotheses (below): the quantum numbers and the
breaking are checked in Section 6 of the hypotheses document (Rust: ALL
PASS (55 of 55); Python: ALL PASS (14 of 14); Wolfram, in exact
arithmetic: ALL PASS (13 of 13)), with the measured inputs of the 2026
edition of the Review of Particle Physics; and the mean-field superfluid
of the three fields is computed in the manuscript
\texttt{higgs\_superfluid.md} (\texttt{.tex}, \texttt{.pdf}), with the
first version as the comparison: in all five methods the ground state is
the custodial Higgs superfluid, which keeps the electric charge and the
custodial symmetry and is stable (Rust: ALL PASS (25 of 25); Python: ALL
PASS (8 of 8); Wolfram: ALL PASS (5 of 5)). The scale, the masses of the
gauge bosons and the amplitude mode are computed in the manuscript
\texttt{higgs\_masses.md} (\texttt{.tex}, \texttt{.pdf}): the photon is
massless and \(\rho_{\rm SH}=1\) in the second version, but the
amplitude mode is \(16\) to \(19\) times too heavy relative to the scale
(Rust: ALL PASS (30 of 30); Python: ALL PASS (10 of 10); Wolfram: ALL
PASS (5 of 5)).

\textbf{The second version of the revised hypotheses} (21 November 2026,
from 19:38 PST). Exploratory computations of stage 4 found that the
superfluid of the first version, with one singlet field, gives
\(\rho\approx1.2\): it pairs only the down component of
\(\Psi_{\rm LL}\) and breaks the custodial symmetry. The author then
instructed to add a second singlet field \(\chi'_{\rm SH}\) with the
hypercharge \(-\tfrac34\). With it the pairs form a bidoublet of
\(SU(2)_L\times SU(2)_{\rm R}\) and the vacuum pairs every constituent.
The document of the hypotheses, its Section 6 and the manuscript
\texttt{higgs\_superfluid.md} state the second version (stages 1 to 3 of
the second version are complete); the manuscript computes the first
version alongside, as the comparison. The first version is kept under
the git tag \texttt{backup/SM-Higgs-one-singlet-2026-11-21}.

\textbf{The determinant interaction} (22 November 2026, from 06:00 PST).
In the second version the particle number is an exact global symmetry
that the condensate breaks, so the superfluid has a fourth Goldstone
mode, a massless physical particle besides the three modes that
\(W^\pm_\mu\) and \(Z_\mu\) absorb. Of the three remedies put to him,
the author chose the determinant interaction
\(-G^{\rm det}_{\rm SH}(\epsilon_{st}\Pi^s_{\rm SH}\Pi'^t_{\rm SH}+{\rm h.c.})\):
it is invariant under \(SU(2)_L\times U(1)_Y\) and the custodial
symmetry, breaks the particle number explicitly, and gives the boson a
mass. The document of the hypotheses states it (H3, consequence 3), and
the task record states the decisions of interpretation (D36) and the
plan (object O19, stage 7). The stages computed before the decision are
its case \(G^{\rm det}_{\rm SH}=0\).

\hypertarget{what-this-folder-is}{%
\section{1. What this folder is}\label{what-this-folder-is}}

This folder of PROJECT spacetimes (repository
\texttt{https://github.com/once-ere/spacetimes.git}) contains the
solution of the problem that the author posed on 20 November 2026: the
gauge-covariant field equations of a self-interacting spinor field on
the \(3+1\) manifold of the Lorentz super-loop, their analytic
solutions, and the condensation of the spinor field into a superfluid
boson singlet state of spin zero, described by three methods
(Hubbard-Stratonovich, Cornwall-Jackiw-Tomboulis and Kohn-Sham), with
the effective Lagrangian of each method, the spontaneous breaking of the
symmetry and the masses of the gauge bosons.

Every document of the folder is complete in itself: every equation,
symbol and definition that it uses is written in it. This README is the
list of contents of the folder.

\hypertarget{state}{%
\section{2. State}\label{state}}

\textbf{20 November 2026.} The folder has been created; the task record
and this README are written. The work paused after that, as the author's
Rule 2 requires after a false statement has been recorded with its cause
(I had stated that the reading of four symbol registries of the earlier
work was complete when two of them had not yet been read completely; the
statement was corrected in the next message). At 18:28 PST the author
instructed to continue and to complete all stages; at 19:25 PST, to push
every item as soon as it is complete. The work continues; this table is
updated with every item.

\textbf{Part I of the manuscript} (items 1 to 5 of the problem: the
manifold M31, the spinor field and theorem D1, the gauge group GEE and
its covariant derivatives, the Lagrangian and the field equations with
their numerical formulation, and four analytic solutions) is complete:
\texttt{superfluid.md}, \texttt{.tex}, \texttt{.pdf}, with the Rust
workspace \texttt{rust/} (program \texttt{sf\_part1}: VERDICT: ALL PASS
(78 of 78)) and the Python programs \texttt{python/sf\_part1.py}
(VERDICT: ALL PASS (34 of 34)) and
\texttt{python/sf\_part1\_figures.py}.

\textbf{Part II of the manuscript} (item 6 of the problem and the three
methods: the pair channels and the gauge bosons that mediate the
pairing, the methods HS, KS and CJT and their numerical engine, the
recursion of the potential to \(V^{(2)}_{\rm SF}\), the ground states,
the first excited states and the gaps, and the gauge bosons in the
superfluid) is complete, with the program \texttt{sf\_part2} (VERDICT:
ALL PASS (44 of 44)) and the Python programs
\texttt{python/sf\_part2.py} (VERDICT: ALL PASS (23 of 23)) and
\texttt{python/sf\_part2\_figures.py}.

\textbf{Part III of the manuscript} (the effective Lagrangians
\(\mathcal L^{\rm CJT}\), \(\mathcal L^{\rm HS}\),
\(\mathcal L^{\rm KS}\), their potential parts (the Mexican hat) and
kinetic parts, the vacua and the breaking of the symmetry, and
\(\mathrm{SFBSv}^X\)) is complete, with the program \texttt{sf\_part3}
(VERDICT: ALL PASS (155 of 155)) and the Python programs
\texttt{python/sf\_part3.py} (VERDICT: ALL PASS (112 of 112)) and
\texttt{python/sf\_part3\_figures.py}. Section 12.1 now also defines the
Landau form, with examples, and states its accuracy, the cause of its
error and why the error matters (Table 12.2, Figure 9); the answer to
the author's questions on this remark is kept in
\texttt{provenance/P4\_remarks.md}. Section 12.4, added at the author's
request of 21 November 2026, computes the free energy of the five
methods at finite temperature (whose value at zero temperature is the
effective potential), the critical temperatures, the gap equation and
the radius of convergence at finite temperature, and answers whether the
order parameter is then smaller and more physically reasonable (Table
12.4, Figure 10); this addendum, and the remark R2 on the mean-field
free energies that it contains, are kept in
\texttt{provenance/P4\_remarks.md} as well.

\textbf{Part IV of the manuscript} (the mass matrix of the gauge bosons
in every vacuum and its diagonalization, the theorem that the massless
gauge bosons are the generators that leave the condensate unchanged, the
surviving group \(SL(2,\mathbb C)\times O(2)\) and its \(U(1)\) with the
massless gauge boson, named the zeugon (not the photon: the spinor
field, the DiracLike spinor, has electric charge 0 and no electroweak
quantum numbers); the dynamics of the superfluid and of its gauge
fields, integrated with CVODE and IDA; two videos) is complete, with the
program \texttt{sf\_part4} (VERDICT: ALL PASS (120 of 120)), the
integrators \texttt{sfsundials} and their C interface
\texttt{sfsundials\_c}, and the Python programs
\texttt{python/sf\_part4.py} (VERDICT: ALL PASS (20 of 20)) and
\texttt{python/sf\_part4\_figures.py}.

\textbf{The Mathematica notebook \texttt{superfluid.nb}} (object O8) is
written part by part. This version contains the introduction, the
register of the checks, Part I, Part II, the first section of Part III
(the Landau form: its definition, examples, accuracy, the cause of its
error and why the error matters), and the second section of Part III
(Section III.T: the free energy of the five methods at finite
temperature, the critical temperatures, the size of the order parameter,
and the remark that the free energies are mean-field free energies), and
the third section of Part III (Section III.G: the gradient coefficient
at the critical temperature, the Ginzburg number and the width of the
critical region, and the superfluid as the Higgs sector of its own gauge
group): 181 checks, each repeated by the Rust and the Python programs
(verdict of the evaluation in the front end of Wolfram: ALL PASS (181 of
181), 0 messages). Since 22 November 2026 it ends with Part R, the
revised hypotheses: the fields and the interaction of the second
version, the 28 checks of the Wolfram scripts of stages 2, 3, 4 and 6 of
the revision with their explanations and their counterparts in Rust and
Python, and the remarks R8 to R12 with the computations they rest on;
the checks above include those of Part R. Its provenance is
\texttt{provenance/P3\_notebook.md}. The rest of Part III and Part IV of
the notebook for the earlier hypotheses are not written.

\textbf{The article for Physical Review D} (folder \texttt{paper\_prd/};
part three of stage 6, 22 November 2026): the Regular Article ``Higgs as
a Composite, and Quantum Gravity'', rewritten for the second version of
the revised hypotheses, with the first version as the comparison
(decision D35 of the task record). It answers, under the revised
hypotheses, the questions of the author's request of 21 November 2026,
11:29 PST: whether the superfluid of each method can play the role of
the Higgs particle, and whether the Higgs is then a composite particle.
The superfluid has the quantum numbers of two Higgs doublets; in all
five methods it breaks \(SU(2)_L\times U(1)_Y\) to the \(U(1)\) of the
electric charge, keeps the photon massless and gives \(\rho=1\), so that
with \(v_{\rm SH}=246.22\) GeV the \(W\) and \(Z\) bosons have their
measured masses at tree level; but its amplitude mode, a composite
scalar, is \(16\) to \(19\) times too heavy relative to the scale, so
within the mean field it is not the Higgs boson. The width of the
critical region, the first question of that request, was answered for
the earlier hypotheses only; the revision has no stage at finite
temperature. With the article come its Supplemental Material, the cover
letter, the information for the submission server of the APS and the
folder of the files to upload, all built by
\texttt{paper\_prd/scripts/build\_paper.py} from the fragments of the
manuscripts of the revision, the data files and the reports of the
checks. The gravitational sector of the title is outlined, not computed.
The submission itself is the author's. The article of the earlier
hypotheses stays unchanged in the folder
\texttt{Revision\_HS\_CJT\_and\_Kohn-Sham/paper\_prd}.

\textbf{Stage 2 of the revision of the hypotheses} (first version 21
November 2026, 17:30 to 18:15 PST; second version, with the two
singlets, 19:50 to 20:35 PST). The program \texttt{sh\_quantum} of the
crate \texttt{rust/shquantum} (ALL PASS (55 of 55)), its second
computation \texttt{python/sh\_quantum.py} (ALL PASS (14 of 14)) and the
exact computation \texttt{wolfram/source/sh\_quantum.wls} (ALL PASS (13
of 13), run in Wolfram 15.0.1) check the quantum numbers of the three
fields and of all their spin-zero bilinears, the custodial symmetry and
the pattern of the breaking. The pairs of the doublet with the first
singlet carry exactly the generators of the Higgs doublet; those with
the second singlet carry those of a doublet whose conjugate is a second
Higgs doublet; together they form a bidoublet of
\(SU(2)_L\times SU(2)_{\rm R}\). The only neutral pairs are those of the
vacuum of the hypotheses. That vacuum keeps the electric charge, and
with equal lengths of its two condensates also the custodial
\(SU(2)_{\rm V}\); it breaks three generators of
\(SU(2)_L\times U(1)_Y\) and the particle number. A kinetic term of the
pair fields with one coefficient each gives the masses of the \(W\) and
the \(Z\) of the Standard Model, the massless photon and \(\rho=1\). The
measured inputs are those of the 2026 edition of the Review of Particle
Physics (\texttt{data/sh\_inputs.csv}).

\textbf{Stage 3 of the revision of the hypotheses} (first version 21
November 2026, 18:15 to 19:22 PST; second version, with the two
singlets, 20:35 to 21:50 PST). The manuscript
\texttt{higgs\_superfluid.md} (\texttt{.tex}, \texttt{.pdf}) computes
the uniform ground states of the three fields by the methods HS (three
parameter sets), KS and CJT, and those of the first version as the
comparison, with the program \texttt{sh\_superfluid} of the crate
\texttt{rust/shsuperfluid} (ALL PASS (25 of 25)), built on the engine
\texttt{uniform\_n} of the library \texttt{sfcore} for several Dirac
fields; the Python program \texttt{python/sh\_superfluid.py} (ALL PASS
(8 of 8)); and the exact computation
\texttt{wolfram/source/sh\_superfluid.wls} (ALL PASS (5 of 5)). The
normal vacuum is unstable only towards the doublet-singlet pairs, with
the critical coupling
\(G^{\rm crit}_{\rm SH}=4\pi^2/I^{(1)}_{\rm SF}(m_{\rm LL})\) derived in
closed form. With the pairing terms, the ground state of every method is
the custodial Higgs superfluid: both pair fields condense with equal
lengths, the down particles paired with \(\chi_{\rm SH}\) and the up
particles with \(\chi'_{\rm SH}\). It keeps exactly the electric charge,
the custodial \(SU(2)_{\rm V}\) and the relative number, has four
Goldstone modes (three for \(W^\pm_\mu\) and \(Z_\mu\), one of the
particle number), and is stable also against charged condensates. Every
constituent is paired, and the up and the down components are equivalent
in it. The condensate of one doublet, the ground state of the first
version, is an unstable stationary state of the second. The remarks R8
to R10 of \texttt{provenance/P4\_remarks.md} explain what forms the
condensate (with the correction of R8), the charges of the lightest
quasi-particles in both versions, and an artifact of the sharp cutoff in
the current response of the vacuum.

\textbf{Stage 4 of the revision of the hypotheses} (21 November 2026,
from 21:50 PST; the former stages 4 and 5, merged). The manuscript
\texttt{higgs\_masses.md} (\texttt{.tex}, \texttt{.pdf}) computes, for
the five methods and both versions, the effective potential of the pair
fields along the direction of the vacuum, the susceptibilities of the
superfluid to the chemical potentials of the eight generators, the scale
\(v_{\rm SH}\) and the masses of the gauge bosons, the parameter
\(\rho_{\rm SH}\), and the amplitude mode by the effective Lagrangian
and by the linear response at a frequency, with a scan of the coupling
and the mass. The programs: \texttt{sh\_higgs} of the crate
\texttt{rust/shhiggs} (ALL PASS (30 of 30)), built on the module
\texttt{rpa\_n} of the library \texttt{sfcore};
\texttt{python/sh\_higgs.py} (ALL PASS (10 of 10)); and
\texttt{wolfram/source/sh\_higgs.wls} (ALL PASS (5 of 5)). The
superfluid is a minimum of its effective potential; the photon is
massless; in the second version \(\rho_{\rm SH}=1\) and four generators
are broken. The top of the effective potential is the normal vacuum,
except in KS and CJT of the second version, where it is the scalar
condensate of stage 3, below zero; the first version of the manuscript
had said that the top is the normal vacuum in every method, and the
correction is recorded in the task record (Section 4.4, entry E60) and
in Section 2.3 of the manuscript. With \(v_{\rm SH}=246.22\) GeV the
cutoff is \(2207\) to \(2413\) GeV, and the amplitude mode is \(16\) to
\(19\) times heavier relative to the scale than the Higgs boson of the
Standard Model. In the pairing-only methods the linear response has the
closed form of the theory of Bardeen, Cooper and Schrieffer, which puts
the amplitude mode exactly at the pair self-energy \(d^0_{\rm SH}\). In
the second version the amplitude mode comes with a custodial triplet of
massive scalars at the same frequency: the second version has the
scalars of two Higgs doublets.

\textbf{Stage 6 of the revision of the hypotheses, part one} (22
November 2026): the dynamics in time. The manuscript
\texttt{higgs\_dynamics.md} (\texttt{.tex}, \texttt{.pdf}) follows the
superfluid by the time-dependent mean field of each of the five methods,
integrated with CVODE: the ground states stay at rest, the energy, the
charges and the orthonormality of the quasi-particle vectors are
conserved, a small displacement of the pair fields oscillates at the
frequencies of the modes, and a chemical potential of the particle
number turns the phase of the pair fields at twice its value (the
Josephson relation). The programs: \texttt{sh\_dynamics} of the crate
\texttt{rust/shdynamics} (ALL PASS (12 of 12));
\texttt{python/sh\_dynamics.py} (ALL PASS (7 of 7)) with the figures by
\texttt{python/sh\_dynamics\_figures.py}; and
\texttt{wolfram/source/sh\_dynamics.wls} (ALL PASS (5 of 5)). In HS (0,
1) and HS (0, 0) the amplitude mode and the custodial triplet oscillate
at \(d^0_{\rm SH}\), undamped, as stage 4 found. In HS (1, 1) the
amplitude mode is bound below the energy that breaks a pair and the
triplet lies 2.2 per cent above it; in KS and CJT the amplitude mode
lies above the energy that breaks a pair and decays. The notebook (part
two of stage 6) is complete: Part R of \texttt{superfluid.nb} holds the
checks and the remarks of the revision. The article (part three) is
complete (above).

\hypertarget{contents}{%
\section{3. Contents}\label{contents}}

\begin{longtable}[]{@{}
  >{\raggedright\arraybackslash}p{(\columnwidth - 4\tabcolsep) * \real{0.3333}}
  >{\raggedright\arraybackslash}p{(\columnwidth - 4\tabcolsep) * \real{0.3333}}
  >{\raggedright\arraybackslash}p{(\columnwidth - 4\tabcolsep) * \real{0.3333}}@{}}
\toprule\noalign{}
\begin{minipage}[b]{\linewidth}\raggedright
file or folder
\end{minipage} & \begin{minipage}[b]{\linewidth}\raggedright
what it is
\end{minipage} & \begin{minipage}[b]{\linewidth}\raggedright
state
\end{minipage} \\
\midrule\noalign{}
\endhead
\bottomrule\noalign{}
\endlastfoot
\texttt{README.md}, \texttt{.tex}, \texttt{.pdf} & this list of contents
& written \\
\texttt{provenance/P0\_task\_record.md}, \texttt{.tex}, \texttt{.pdf} &
the task (paraphrased), how the author's rules are kept, the record of a
false statement and its correction, the decisions of interpretation, the
plan, the log of errors and changes & written \\
\texttt{provenance/H0\_hypotheses.md}, \texttt{.tex}, \texttt{.pdf} &
the revised hypotheses of this folder: the fields, the gauge group with
the hypercharge, the pairing term of the doublet-singlet channel, the
quantum numbers of the superfluid and why they are those of the Higgs
doublet, the scale, and the consequences that follow at once & written
(current) \\
\texttt{higgs\_dynamics.md}, \texttt{.tex}, \texttt{.pdf} & the
manuscript of stage 6 of the revision, part one: the time-dependent mean
field of the five methods, its conservation laws, the amplitude mode and
the custodial triplet in the time domain, the Josephson relation, a
large displacement, the checks & written (current, second version) \\
\texttt{higgs\_masses.md}, \texttt{.tex}, \texttt{.pdf} & the manuscript
of stage 4 of the revision: the effective potential, the
susceptibilities to the chemical potentials of the generators, the scale
and the masses of the gauge bosons, the parameter \(\rho\), the
amplitude mode by the effective Lagrangian and by the linear response,
the scan of the coupling and the mass, the checks & written (current,
second version) \\
\texttt{higgs\_superfluid.md}, \texttt{.tex}, \texttt{.pdf} & the
manuscript of stage 3 of the revision: the mean-field superfluid of the
three fields of the second version of the revised hypotheses by the
methods HS, KS and CJT, with the first version as the comparison (the
Hamiltonian density and the potentials, what forms the condensate, the
uniform states and the three methods, the static response of the normal
vacuum with the critical coupling in closed form, the stationary states
and the custodial Higgs ground states, their symmetry and stability, the
instability of the condensate of one doublet, the quasi-particles and
the custodial symmetry, the checks) & written (current, second
version) \\
\texttt{provenance/P1\_master\_symbols.md}, \texttt{.tex}, \texttt{.pdf}
& the MASTER provenance of the symbols of PROJECT spacetimes for this
solution: the convention of the symbols (suffix LL: taken unchanged from
the earlier work on the super-loop; suffix SF: introduced by this
solution), the seven registries of the project read by a program, the
identifications and remedies file by file, the collisions and their
resolution, the machine checks, and the table of all symbols & written;
grows with the solution \\
\texttt{provenance/symbols\_sf.csv} & the registry of the symbols of the
solution, written by \texttt{code/scripts/write\_symbol\_registry.py}
from \texttt{code/scripts/symbol\_entries.py} & written \\
\texttt{superfluid.md}, \texttt{.tex}, \texttt{.pdf} & the manuscript of
the solution & parts I to IV written \\
\texttt{superfluid.nb}, \texttt{superfluid\_notebook.pdf} & the
Mathematica notebook of the whole solution, evaluated in the front end
of Wolfram, and its printed copy & introduction, Parts I and II and the
sections of Part III on the Landau form (III.L), on the free energy at
finite temperature (III.T) and on the gradient coefficient, the Ginzburg
number and the Higgs sector (III.G) written, and Part R (the revised
hypotheses); the rest of Parts III and IV for the earlier hypotheses not
written \\
\texttt{provenance/P3\_notebook.md}, \texttt{.tex}, \texttt{.pdf} & the
provenance of the notebook: its task, files, build and evaluation, the
checks with their counterparts in Rust and Python, the symbols of its
code, the corrections & written for Parts I to III and Part R; grows
with the notebook \\
\texttt{provenance/P4\_remarks.md}, \texttt{.tex}, \texttt{.pdf} & the
remarks on the results (an approximation, a fit, an accuracy, a
limitation): for each, what it means, why it holds, why it matters, the
corrections of the first answer, and the provenance; first remark R1:
the accuracy of the Landau form, with its definition and examples; the
addendum A1 (the free energy of the five methods at finite temperature,
and whether the order parameter is smaller and more physically
reasonable) and remark R2 (the free energies are mean-field free
energies); the addendum A2 (the gradient coefficient at the critical
temperature and the width of the critical region; the superfluid as a
Higgs sector and the compositeness of the Higgs) with its remarks; for
the revised hypotheses the remarks R8 to R10 of stage 3 (what forms the
condensate, with the correction of R8; the charges of the lightest
quasi-particles in the first and the second version; the current
response of the vacuum with the sharp cutoff) & written for R1, A1, R2,
A2 and R8 to R10; grows with every such remark \\
\texttt{provenance/} (further documents) & the provenance of the code,
the computed values and the checks & planned \\
\texttt{code/scripts/} & the document pipeline:
\texttt{assemble\_docs.py} (writes a document from its source in
\texttt{docs\_src/} and the generated fragments),
\texttt{build\_docs.sh} (the \texttt{.tex} and the \texttt{.pdf} of
every document, with a check of the layout of every page),
\texttt{doc\_header.tex}, \texttt{keep\_tables.py},
\texttt{check\_layout.py}, \texttt{check\_pages.py}, and
\texttt{check\_self\_contained.py} (does a document send its reader
elsewhere?); the registry of the symbols: \texttt{symbol\_entries.py}
(the entries), \texttt{write\_symbol\_registry.py} (reads the seven
registries of the project, writes the registry of the solution and the
tables of the MASTER provenance) and \texttt{check\_symbols.py} (the
checks of the symbols) & written \\
\texttt{docs\_src/} & the hand-edited sources of the documents &
written \\
\texttt{rust/} & the Rust workspace: the library \texttt{sfcore}
(bi-quaternions and bi-octonions, Dirac matrices, automatic
differentiation to second order, Fock-space operators, eigenvalues of
block-tridiagonal matrices, quadrature, reports; for part II the uniform
quasi-free states of the three methods, the linear response, the
eigenvalues of general complex matrices and the exchange of gauge
bosons; for part III the states with a fixed pair field and the thermal
states at a temperature \(T\); for part IV the response at complex
frequencies), the program \texttt{sf\_part1} (the checks and analytic
solutions of part I), the program \texttt{sf\_part2} (part II), the
program \texttt{sf\_part3} (part III), the program \texttt{sf\_part4}
(part IV), the integrators CVODE and IDA of the vendored pure-Rust
SUNDIALS 7.8.0 with a plain interface (\texttt{sfsundials}), and their C
interface for Python, the shared library \texttt{libsfsundials\_c}
(\texttt{sfsundials\_c}); for the revision of the hypotheses the program
\texttt{sh\_quantum} (crate \texttt{shquantum}: the quantum numbers of
the fields and of their bilinears, the breaking, the masses of the gauge
fields from the pair field), the module \texttt{uniform\_n} of
\texttt{sfcore} (the uniform quasi-free states of several Dirac fields
and their static response) the program \texttt{sh\_superfluid} (crate
\texttt{shsuperfluid}: the mean-field superfluid of the three fields),
the module \texttt{rpa\_n} of \texttt{sfcore} (the linear response at a
frequency of the states of several Dirac fields) the program
\texttt{sh\_higgs} (crate \texttt{shhiggs}: the effective potential, the
susceptibilities, the masses of the gauge bosons and the amplitude mode)
and the program \texttt{sh\_dynamics} (crate \texttt{shdynamics}: the
time-dependent mean field of the five methods, integrated with CVODE) &
parts I to IV written; stages 2 to 4 and the first part of stage 6 of
the revision written \\
\texttt{python/} & \texttt{sfcommon.py} (common definitions),
\texttt{sfsundials.py} (CVODE and IDA through the C interface),
\texttt{sf\_part1.py}, \texttt{sf\_part2.py}, \texttt{sf\_part3.py},
\texttt{sf\_part4.py} (the second, independent computations of parts I
to IV), \texttt{sf\_part1\_figures.py}, \texttt{sf\_part2\_figures.py},
\texttt{sf\_part3\_figures.py}, \texttt{sf\_part4\_figures.py} (the
figures, and the videos of part IV); \texttt{sh\_quantum.py} (the second
computation of \texttt{sh\_quantum}); \texttt{sh\_superfluid.py} (the
second computation of \texttt{sh\_superfluid}, with its own engine and
the closed forms of the pairing-only ground state and of the current
response); \texttt{sh\_higgs.py} (the second computation of
\texttt{sh\_higgs}: the susceptibilities from the energy, the curvature
and the amplitude mode of the pairing-only method);
\texttt{sh\_dynamics.py} (the second computation of
\texttt{sh\_dynamics}, with its own right side of the equations of
motion) and \texttt{sh\_dynamics\_figures.py} (its figures);
\texttt{sh\_higgs\_figures.py} (the figure of the effective potential of
stage 4) and \texttt{sh\_article\_figures.py} (the three figures of the
article) & parts I to IV written; stages 2 to 4 and stage 6 of the
revision written \\
\texttt{wolfram/} & the Wolfram-language parts of the notebook: the
sources \texttt{wolfram/source/superfluid.source.wl},
\texttt{part1.source.wl}, \texttt{part2.source.wl},
\texttt{part3\_landau.source.wl}, \texttt{part3\_thermal.source.wl},
\texttt{part3\_ginzburg.source.wl} and \texttt{part\_revision.source.wl}
(Part R), and the data cells \texttt{data\_part2.wl},
\texttt{data\_part3.wl}, \texttt{data\_part3\_thermal.wl},
\texttt{data\_part3\_ginzburg.wl} and \texttt{data\_partR.wl}; in
\texttt{wolfram/tools/} the builder of the notebook from its source
(\texttt{build\_notebook.wls}), the reader of formulas that turns the
TeX of the texts into two-dimensional formulas
(\texttt{sf\_formula.wl}), the evaluation of the notebook in the front
end of Wolfram (\texttt{guirun.wl} and the three calls
\texttt{gui\_start.wl}, \texttt{gui\_go.wl}, \texttt{gui\_status.wl}),
the dry run (\texttt{make\_dryrun.py}), and the writer of the data cells
(\texttt{make\_data\_cells.py}); the exact computations of the checks of
\texttt{sh\_quantum} and \texttt{sh\_superfluid},
\texttt{wolfram/source/sh\_quantum.wls} and
\texttt{wolfram/source/sh\_superfluid.wls}, of \texttt{sh\_higgs},
\texttt{wolfram/source/sh\_higgs.wls}, and of \texttt{sh\_dynamics},
\texttt{wolfram/source/sh\_dynamics.wls} & Parts I and II, the sections
III.L, III.T and III.G of Part III and Part R written; stages 2 to 4 and
the first two parts of stage 6 of the revision written \\
\texttt{data/} & the reports of the checks
(\texttt{hck\_part1\_checks.txt},
\texttt{hck\_part1\_checks\_python.txt},
\texttt{hck\_part2\_checks.txt},
\texttt{hck\_part2\_checks\_python.txt},
\texttt{hck\_part3\_checks.txt},
\texttt{hck\_part3\_checks\_python.txt},
\texttt{hck\_part4\_checks.txt},
\texttt{hck\_part4\_checks\_python.txt},
\texttt{hck\_symbol\_checks.txt},
\texttt{hck\_self\_contained\_checks.txt}) and the data of parts I to IV
(\texttt{part1\_*.csv}, \texttt{part2\_*.csv}, \texttt{part3\_*.csv},
\texttt{part4\_*.csv}); for the revision of the hypotheses the measured
inputs \texttt{sh\_inputs.csv}, the reports
\texttt{hck\_sh\_quantum\_checks.txt},
\texttt{hck\_sh\_quantum\_checks\_python.txt},
\texttt{hck\_sh\_quantum\_checks\_wolfram.txt} and the data
\texttt{sh\_quantum\_numbers.csv}, \texttt{sh\_channels.csv},
\texttt{sh\_custodial\_channels.csv}, \texttt{sh\_couplings.csv},
\texttt{sh\_mass\_matrix.csv}, \texttt{sh\_quantum\_values.csv}; of
stage 3 the reports \texttt{hck\_sh\_superfluid\_checks.txt},
\texttt{hck\_sh\_superfluid\_checks\_python.txt},
\texttt{hck\_sh\_superfluid\_checks\_wolfram.txt} and the data
\texttt{sh\_sf\_normal\_vacuum.csv}, \texttt{sh\_sf\_states.csv},
\texttt{sh\_sf\_ground.csv}, \texttt{sh\_sf\_values.csv}; of stage 4 the
reports \texttt{hck\_sh\_higgs\_checks.txt},
\texttt{hck\_sh\_higgs\_checks\_python.txt},
\texttt{hck\_sh\_higgs\_checks\_wolfram.txt} and the data
\texttt{sh\_hg\_potential.csv},
\texttt{sh\_hg\_potential\_second\_branch.csv},
\texttt{sh\_hg\_minimum.csv}, \texttt{sh\_hg\_susceptibility.csv},
\texttt{sh\_hg\_masses.csv}, \texttt{sh\_hg\_scan.csv},
\texttt{sh\_hg\_values.csv}; of stage 6, part one, the reports
\texttt{hck\_sh\_dynamics\_checks.txt},
\texttt{hck\_sh\_dynamics\_checks\_python.txt},
\texttt{hck\_sh\_dynamics\_checks\_wolfram.txt} and the data
\texttt{sh\_dyn\_runs.csv}, \texttt{sh\_dyn\_signals.csv},
\texttt{sh\_dyn\_spectra.csv}, \texttt{sh\_dyn\_values.csv} & parts I to
IV written; stages 2 to 4 and the first part of stage 6 of the revision
written \\
\texttt{figures/}, \texttt{videos/} & the figures
(\texttt{fig\_part1\_*}, \texttt{fig\_part2\_*}, \texttt{fig\_part3\_*},
\texttt{fig\_part4\_*}, and of the revision
\texttt{fig\_sh\_hg\_potential}, \texttt{fig\_sh\_dyn\_*} and the
figures of the article \texttt{fig\_prd\_sh\_potential},
\texttt{fig\_prd\_sh\_scan}, \texttt{fig\_prd\_sh\_modes}, each
\texttt{.pdf} and \texttt{.png}) and the videos
(\texttt{video\_part4\_v1\_condensate\_growth.mp4},
\texttt{video\_part4\_v2\_gauge\_fields.mp4}) & written \\
\texttt{code/scripts/write\_part1\_fragments.py},
\texttt{write\_part2\_fragments.py},
\texttt{write\_part3\_fragments.py}, \texttt{write\_part4\_fragments.py}
& turn the reports and values of parts I to IV into the tables and
numbers of the manuscript & written \\
\texttt{code/scripts/write\_sh\_fragments.py} & turns the reports and
data of stage 2 of the revision into the tables and numbers of the
hypotheses document (\texttt{provenance/generated/sh\_quantum.md}) &
written \\
\texttt{code/scripts/write\_sh\_superfluid\_fragments.py} & turns the
reports and data of stage 3 of the revision into the tables and numbers
of \texttt{higgs\_superfluid.md} and of the remarks R8 to R10
(\texttt{provenance/generated/sh\_superfluid.md}) & written \\
\texttt{code/scripts/write\_sh\_dynamics\_fragments.py} & turns the
reports and data of stage 6, part one, into the tables and numbers of
\texttt{higgs\_dynamics.md}
(\texttt{provenance/generated/sh\_dynamics.md}) & written \\
\texttt{code/scripts/write\_sh\_higgs\_fragments.py} & turns the reports
and data of stage 4 of the revision into the tables and numbers of
\texttt{higgs\_masses.md} (\texttt{provenance/generated/sh\_higgs.md}) &
written \\
\texttt{code/scripts/write\_bound\_fragments.py} & writes every bound of
agreement stated in the manuscript and in
\texttt{provenance/P4\_remarks.md} from the reports of the checks (the
largest deviation, rounded up) into
\texttt{provenance/generated/bounds.md} & written \\
\texttt{paper\_prd/} & the Regular Article ``Higgs as a Composite, and
Quantum Gravity'' for Physical Review D
(\texttt{higgs\_composite\_qg.pdf}, \texttt{.tex}, \texttt{.md}), its
Supplemental Material (\texttt{higgs\_composite\_qg\_SM.*}), the cover
letter (\texttt{cover\_letter.*}), the information for the submission
server (\texttt{submission\_info.*}), the folder \texttt{submission/}
with exactly the files to upload, the templates \texttt{src/}, the build
\texttt{scripts/build\_paper.py} with its record
\texttt{generated/numbers.csv} and \texttt{generated/build\_report.txt},
and its README (\texttt{paper\_prd/README.md}, \texttt{.tex},
\texttt{.pdf}) & written; not submitted (the submission is the
author's) \\
\end{longtable}

\hypertarget{how-to-build-the-documents}{%
\section{4. How to build the
documents}\label{how-to-build-the-documents}}

The programs are run first, then the documents are built, from the
folder \texttt{Revision\_SM\_Higgs}:

\begin{Shaded}
\begin{Highlighting}[]
\BuiltInTok{cd} \StringTok{"/Users/nsh/Developer/github/spacetimes/Revision\_SM\_Higgs"}
\KeywordTok{(}\BuiltInTok{cd}\NormalTok{ rust }\KeywordTok{\&\&} \ExtensionTok{cargo}\NormalTok{ build }\AttributeTok{{-}{-}release} \KeywordTok{\&\&} \ExtensionTok{./target/release/sf\_part1}\KeywordTok{)}   \CommentTok{\# part I: checks and data}
\ExtensionTok{python3}\NormalTok{ python/sf\_part1.py                    }\CommentTok{\# part I: the Python computation of the checks}
\ExtensionTok{python3}\NormalTok{ python/sf\_part1\_figures.py            }\CommentTok{\# part I: the figures}
\ExtensionTok{python3}\NormalTok{ code/scripts/write\_part1\_fragments.py }\CommentTok{\# the tables and numbers of part I for the manuscript}
\KeywordTok{(}\BuiltInTok{cd}\NormalTok{ rust }\KeywordTok{\&\&} \ExtensionTok{./target/release/sf\_part2}\KeywordTok{)}        \CommentTok{\# part II: checks and data (about nine minutes)}
\ExtensionTok{python3}\NormalTok{ python/sf\_part2.py                    }\CommentTok{\# part II: the Python computation (after sf\_part2)}
\ExtensionTok{python3}\NormalTok{ python/sf\_part2\_figures.py            }\CommentTok{\# part II: the figures}
\ExtensionTok{python3}\NormalTok{ code/scripts/write\_part2\_fragments.py }\CommentTok{\# the tables and numbers of part II for the manuscript}
\KeywordTok{(}\BuiltInTok{cd}\NormalTok{ rust }\KeywordTok{\&\&} \ExtensionTok{./target/release/sf\_part3}\KeywordTok{)}        \CommentTok{\# part III: checks and data (after sf\_part2)}
\ExtensionTok{python3}\NormalTok{ python/sf\_part3.py                    }\CommentTok{\# part III: the Python computation}
\ExtensionTok{python3}\NormalTok{ python/sf\_part3\_figures.py            }\CommentTok{\# part III: the figures}
\ExtensionTok{python3}\NormalTok{ code/scripts/write\_part3\_fragments.py }\CommentTok{\# the tables and numbers of part III for the manuscript}
\KeywordTok{(}\BuiltInTok{cd}\NormalTok{ rust }\KeywordTok{\&\&} \ExtensionTok{./target/release/sf\_part4}\KeywordTok{)}        \CommentTok{\# part IV: checks and data (after sf\_part3; about twenty minutes)}
\ExtensionTok{python3}\NormalTok{ python/sf\_part4.py                    }\CommentTok{\# part IV: the Python computation (uses the library of cargo build)}
\ExtensionTok{python3}\NormalTok{ python/sf\_part4\_figures.py            }\CommentTok{\# part IV: the figures and the two videos (needs ffmpeg)}
\ExtensionTok{python3}\NormalTok{ code/scripts/write\_part4\_fragments.py }\CommentTok{\# the tables and numbers of part IV for the manuscript}
\ExtensionTok{python3}\NormalTok{ code/scripts/write\_bound\_fragments.py }\CommentTok{\# the bounds of agreement of the documents (after the programs of parts I to IV)}
\KeywordTok{(}\BuiltInTok{cd}\NormalTok{ rust }\KeywordTok{\&\&} \ExtensionTok{./target/release/sh\_quantum}\KeywordTok{)}      \CommentTok{\# revision, stage 2: the quantum numbers and the breaking (under a second)}
\ExtensionTok{python3}\NormalTok{ python/sh\_quantum.py                  }\CommentTok{\# revision, stage 2: the Python computation (after sh\_quantum)}
\ExtensionTok{wolframscript} \AttributeTok{{-}file}\NormalTok{ wolfram/source/sh\_quantum.wls data/sh\_channels.csv data/sh\_custodial\_channels.csv }\DataTypeTok{\textbackslash{}}
\NormalTok{    data/hck\_sh\_quantum\_checks\_wolfram.txt}
\ExtensionTok{python3}\NormalTok{ code/scripts/write\_sh\_fragments.py    }\CommentTok{\# the tables and numbers of stage 2 for the hypotheses document}
\KeywordTok{(}\BuiltInTok{cd}\NormalTok{ rust }\KeywordTok{\&\&} \ExtensionTok{./target/release/sh\_superfluid}\KeywordTok{)}   \CommentTok{\# revision, stage 3: the superfluid (about fifteen minutes on two threads)}
\KeywordTok{(}\BuiltInTok{cd}\NormalTok{ python }\KeywordTok{\&\&} \ExtensionTok{python3}\NormalTok{ sh\_superfluid.py}\KeywordTok{)}       \CommentTok{\# revision, stage 3: the Python computation (after sh\_superfluid)}
\ExtensionTok{wolframscript} \AttributeTok{{-}file}\NormalTok{ wolfram/source/sh\_superfluid.wls data/sh\_sf\_normal\_vacuum.csv data/sh\_sf\_ground.csv }\DataTypeTok{\textbackslash{}}
\NormalTok{    data/hck\_sh\_superfluid\_checks\_wolfram.txt}
\ExtensionTok{python3}\NormalTok{ code/scripts/write\_sh\_superfluid\_fragments.py }\CommentTok{\# the tables and numbers of stage 3}
\KeywordTok{(}\BuiltInTok{cd}\NormalTok{ rust }\KeywordTok{\&\&} \ExtensionTok{./target/release/sh\_higgs}\KeywordTok{)}        \CommentTok{\# revision, stage 4: the masses (about an hour on two threads)}
\KeywordTok{(}\BuiltInTok{cd}\NormalTok{ python }\KeywordTok{\&\&} \ExtensionTok{python3}\NormalTok{ sh\_higgs.py}\KeywordTok{)}            \CommentTok{\# revision, stage 4: the Python computation (after sh\_higgs)}
\ExtensionTok{wolframscript} \AttributeTok{{-}file}\NormalTok{ wolfram/source/sh\_higgs.wls data/sh\_hg\_minimum.csv data/sh\_hg\_masses.csv data/sh\_hg\_values.csv }\DataTypeTok{\textbackslash{}}
\NormalTok{    data/hck\_sh\_higgs\_checks\_wolfram.txt}
\ExtensionTok{python3}\NormalTok{ python/sh\_higgs\_figures.py            }\CommentTok{\# the figure of stage 4 (the effective potential)}
\ExtensionTok{python3}\NormalTok{ code/scripts/write\_sh\_higgs\_fragments.py }\CommentTok{\# the tables and numbers of stage 4}
\KeywordTok{(}\BuiltInTok{cd}\NormalTok{ rust }\KeywordTok{\&\&} \ExtensionTok{./target/release/sh\_dynamics}\KeywordTok{)}     \CommentTok{\# revision, stage 6, part one: the dynamics (about forty minutes)}
\KeywordTok{(}\BuiltInTok{cd}\NormalTok{ python }\KeywordTok{\&\&} \ExtensionTok{python3}\NormalTok{ sh\_dynamics.py}\KeywordTok{)}         \CommentTok{\# revision, stage 6, part one: the Python computation}
\ExtensionTok{python3}\NormalTok{ python/sh\_dynamics\_figures.py         }\CommentTok{\# the figures of the dynamics}
\ExtensionTok{wolframscript} \AttributeTok{{-}file}\NormalTok{ wolfram/source/sh\_dynamics.wls data/sh\_dyn\_runs.csv data/sh\_dyn\_signals.csv }\DataTypeTok{\textbackslash{}}
\NormalTok{    data/hck\_sh\_dynamics\_checks\_wolfram.txt}
\ExtensionTok{python3}\NormalTok{ code/scripts/write\_sh\_dynamics\_fragments.py }\CommentTok{\# the tables and numbers of stage 6, part one}
\ExtensionTok{python3}\NormalTok{ code/scripts/write\_notebook\_fragments.py }\CommentTok{\# the tables of the provenance of the notebook (from data/hck\_notebook\_log.txt)}
\ExtensionTok{python3}\NormalTok{ code/scripts/write\_symbol\_registry.py }\CommentTok{\# the registry of the symbols and its tables}
\ExtensionTok{python3}\NormalTok{ code/scripts/check\_symbols.py         }\CommentTok{\# writes data/hck\_symbol\_checks.txt}
\ExtensionTok{python3}\NormalTok{ code/scripts/assemble\_docs.py         }\CommentTok{\# docs\_src/*.src.md {-}\textgreater{} the .md documents}
\FunctionTok{bash}\NormalTok{ code/scripts/build\_docs.sh               }\CommentTok{\# .md {-}\textgreater{} .tex {-}\textgreater{} .pdf, with the layout check}
\ExtensionTok{python3}\NormalTok{ code/scripts/check\_self\_contained.py  }\CommentTok{\# writes data/hck\_self\_contained\_checks.txt}
\ExtensionTok{python3}\NormalTok{ python/sh\_article\_figures.py          }\CommentTok{\# the three figures of the article}
\ExtensionTok{python3}\NormalTok{ paper\_prd/scripts/build\_paper.py      }\CommentTok{\# the article for Physical Review D, its companions}
\FunctionTok{bash}\NormalTok{ code/scripts/build\_docs.sh paper\_prd/README.md  }\CommentTok{\# the README of paper\_prd/ (named)}
\end{Highlighting}
\end{Shaded}

The Rust programs need the Rust toolchain (\texttt{cargo}; no crate is
downloaded: the integrators are the pure-Rust SUNDIALS 7.8.0 vendored in
PROJECT spacetimes as \texttt{vendor/sundials\_rs-7.8.0}); the Python
programs need Python 3 with numpy, scipy, sympy and matplotlib, and the
videos the program ffmpeg. \texttt{assemble\_docs.py} writes each
document from its source in \texttt{docs\_src/}; \texttt{build\_docs.sh}
needs pandoc (version 2.11 or later), a TeX distribution with XeLaTeX
and latexmk, the DejaVu fonts and the font Latin Modern Math;
\texttt{check\_self\_contained.py} checks that no document refers its
reader to another document. \texttt{build\_paper.py} needs a TeX
installation with REVTeX 4.2 and pdfLaTeX, and pandoc; the default run
of \texttt{build\_docs.sh} leaves the folder \texttt{paper\_prd/} alone,
because the \texttt{.tex} files there are REVTeX sources written by
\texttt{build\_paper.py}.

The notebook is opened in Wolfram (version 15 or later) and evaluated
with Evaluation \textgreater{} Evaluate Notebook. It is made from its
source on a computer with Wolfram by

\begin{Shaded}
\begin{Highlighting}[]
\ExtensionTok{wolframscript} \AttributeTok{{-}file}\NormalTok{ wolfram/tools/build\_notebook.wls wolfram/source/superfluid.source.wl built.nb wolfram/tools/sf\_formula.wl}
\end{Highlighting}
\end{Shaded}

which writes the unevaluated notebook \texttt{built.nb}; evaluating it
gives \texttt{superfluid.nb} (the evaluation of record, with its log
\texttt{data/hck\_notebook\_log.txt}, is made in the front end by
\texttt{wolfram/tools/guirun.wl};
\texttt{python3\ wolfram/tools/make\_dryrun.py\ dry.wls\ wolfram/source/superfluid.source.wl}
and \texttt{wolframscript\ -file\ dry.wls} evaluate the code of all
input cells without the front end).

\end{document}
